\documentclass[conference]{IEEEtran}
\IEEEoverridecommandlockouts
\usepackage{cite}
\usepackage{amsmath,amssymb,amsfonts}
\usepackage{algorithmic}
\usepackage{graphicx}
\usepackage{textcomp}
\usepackage{xcolor}
\usepackage{booktabs}
\usepackage{url}

\def\BibTeX{{\rm B\kern-.05em{\sc i\kern-.025em b}\kern-.08em
    T\kern-.1667em\lower.7ex\hbox{E}\kern-.125emX}}

\begin{document}

\title{A Dual-Mode, Fault-Tolerant Hybrid Architecture for Lexical-Semantic Candidate Ranking in Applicant Tracking Systems}

\author{\IEEEauthorblockN{Anonymous Authors}
\IEEEauthorblockA{\textit{Department of Computer Science and Engineering} \\
\textit{Affiliation Name}\\
City, Country \\
email@institution.edu}
}

\maketitle

\begin{abstract}
Automated Applicant Tracking Systems (ATS) face a critical trilemma between semantic comprehension, exact technical skill verification, and operational resilience. While modern Large Language Models (LLMs) and dense text embeddings capture high-level contextual semantics, they frequently suffer from high inference latency, recurring API costs, susceptibility to external network outages, and hallucinations regarding mandatory hard skills. Conversely, traditional lexical search heuristics (e.g., BM25) ensure strict skill verification and minimal latency, but fail to bridge the semantic vocabulary gap between candidates and job descriptions. 
In this paper, we propose \textbf{Quick-Hire}, an enterprise-grade, dual-mode ATS architecture that unifies dense semantic vector representations, structured lexical coverage, and localized fallback heuristics into a mathematically grounded ranking objective. We formulate a parameterized scoring function $\mathcal{S}(R, J) = \alpha \cdot \cos(\mathbf{e}_R, \mathbf{e}_J) + \beta \cdot \text{Overlap}(S_R, S_J) + \gamma \cdot \mathbb{I}(T_R \sim T_J)$ and evaluate its performance against classic lexical baselines (BM25, Jaccard) and pure dense embedding models across 250 annotated resume-job evaluation pairs. Empirical results demonstrate that our hybrid model achieves an NDCG@5 of \textbf{0.9700} and a Precision@5 of \textbf{0.8000}, significantly outperforming pure dense embeddings (NDCG@5: 0.9318, P@5: 0.7000) and BM25 (P@5: 0.7600), while our zero-cost deterministic offline fallback guarantees 100\% MRR during cloud connectivity disruptions.
\end{abstract}

\begin{IEEEkeywords}
Applicant Tracking Systems, Information Retrieval, Large Language Models, Hybrid Ranking, Dense Embeddings, Fault-Tolerant Architecture, Natural Language Processing.
\end{IEEEkeywords}

\section{Introduction}
The exponential expansion of digital recruitment platforms has led to an influx of hundreds of unstructured resumes per corporate job vacancy. To alleviate the human bottleneck of initial resume triaging, Applicant Tracking Systems (ATS) increasingly rely on automated ranking algorithms. Historically, these systems utilized lexical information retrieval algorithms, such as Term Frequency-Inverse Document Frequency (TF-IDF) and BM25 \cite{robertson2009probabilistic}, to rank candidates based on keyword matches. 

However, lexical methods suffer fundamentally from the \textit{vocabulary mismatch problem}: a candidate describing extensive experience in ``distributed data streaming'' may fail to match a requisition requiring ``Apache Kafka,'' despite having identical functional competency. 

The advent of transformer-based bi-encoders and Large Language Models (LLMs) has introduced dense semantic embeddings \cite{reimers2019sentence, geminiteam2024gemini}, which represent documents as fixed-dimensional vectors in continuous semantic space. While effective at capturing conceptual proximity, pure dense retrieval introduces several critical vulnerabilities into production recruitment workflows:
\begin{enumerate}
    \item \textbf{Loss of Strict Technical Gatekeeping:} High cosine similarity can occur between generalist profiles and specialized roles, inadvertently promoting candidates who lack non-negotiable mandatory skills (e.g., matching a senior C++ kernel programmer to an asynchronous Python role due to shared distributed systems vocabulary).
    \item \textbf{Operational Fragility and High Latency:} Relying purely on cloud-hosted LLM endpoints introduces network latency (often exceeding 1,200\,ms per document) and single points of failure during rate-limiting or provider outages.
    \item \textbf{Opacity and Lack of Recruiter Recourse:} Monolithic neural similarity scores provide recruiters with zero visibility into missing technical competencies.
\end{enumerate}

To overcome these challenges, we introduce \textbf{Quick-Hire}, an open, dual-mode ATS architecture that synergizes dense vector embeddings, structured lexical coverage, and localized fallback engines into an explainable, parameter-tuned ranking system. 

\textbf{Key Contributions of this Work:}
\begin{itemize}
    \item \textbf{Architectural Dual-Mode Resilience:} We design a fault-tolerant ATS system that automatically toggles between cloud LLM inference and a deterministic, zero-latency local fallback engine during API failures.
    \item \textbf{Empirically Tuned Hybrid Objective:} We formalize and evaluate a hybrid ranking metric combining semantic cosine similarity, exact skill coverage, and role-alignment bonuses, optimizing ranking precision over 250 evaluation pairs.
    \item \textbf{Comprehensive Empirical Evaluation:} We benchmark Quick-Hire against classical IR algorithms (BM25, Jaccard) and pure dense embeddings across standard metrics (NDCG@K, Precision@K, MRR, Latency), accompanied by an ablation study.
    \item \textbf{Explainability and Privacy:} We incorporate token-level skill gap extraction (matched vs. missing skills) and PII decoupling to support transparent and equitable hiring decisions.
\end{itemize}

\section{Related Work}

\subsection{Lexical Retrieval in Automated Recruitment}
Early recruitment systems relied on relational database queries and inverted indices powered by BM25 \cite{robertson2009probabilistic}. While computationally lightweight ($\le 0.1$\,ms latency), lexical matching is notoriously brittle against synonymy and polysemy. Recent efforts such as LiJAR \cite{borisyuk2016lijar} and Skill2Vec \cite{li2021skill2vec} demonstrated that keyword matching without semantic enrichment leads to suboptimal candidate discovery for specialized engineering roles.

\subsection{Dense Embeddings and Deep Person-Job Fit}
With the introduction of Siamese transformers \cite{reimers2019sentence, vaswani2017attention}, dense representation learning has become dominant in recommendation systems. Zhu et al. \cite{zhu2018person} pioneered multi-task deep representation learning for person-job fit. However, deep neural rankers frequently suffer from ``hallucinated relevance,'' where semantic distance fails to penalize missing prerequisite qualifications. Our approach mitigates this by enforcing an explicit dual-objective constraint combining dense vectors with exact lexical sets.

\subsection{Fairness and Bias in Algorithmic Hiring}
Algorithmic recruitment has raised significant concerns regarding gender and racial bias encoded within latent text embeddings \cite{raghavan2020mitigating, dwork2012fairness}. To prevent proxy discrimination, modern ATS frameworks must decouple Personally Identifiable Information (PII) from semantic vectorization, a principle integrated into our system design.

\section{System Architecture}
Quick-Hire is implemented as a decoupled, multi-tier system comprising a Spring Boot microservice backend, an Apache PDFBox ingestion layer, an asynchronous LLM orchestration client, and an in-memory transactional database.

\begin{figure}[htbp]
\centering
\fbox{\parbox{0.92\linewidth}{\small
\textbf{Candidate Intake:} PDF Ingestion $\rightarrow$ PII Masking $\rightarrow$ Structured Parsing\\
\hspace*{1.2cm} $\Downarrow$\\
\textbf{Dual-Mode Router:} Is Cloud API Available?\\
\hspace*{0.5cm} $\swarrow$ [Yes] \hspace{2.2cm} $\searrow$ [No]\\
\textbf{Cloud Pipeline:} \hspace{2.0cm} \textbf{Local Fallback Engine:}\\
• LLM JSON Extractor \hspace{1.1cm} • Deterministic Regex / Lexicon\\
• 768-dim Vector Embeddings \hspace{0.4cm} • 768-dim Hash-Projected IDF\\
\hspace*{1.2cm} $\searrow$ \hspace{2.2cm} $\swarrow$\\
\textbf{Hybrid Ranking Engine:} $\mathcal{S}(R, J) = \alpha \cdot \cos + \beta \cdot \text{Overlap} + \gamma \cdot \text{Bonus}$\\
\hspace*{1.2cm} $\Downarrow$\\
\textbf{Ranked Candidate Output with Recruiter Explanations}
}}
\caption{The Quick-Hire Dual-Mode Architecture.}
\label{fig:arch}
\end{figure}

\subsection{Ingestion and Document Parsing}
When a candidate submits a PDF resume, Apache PDFBox extracts raw character streams. To adhere to fair hiring standards, a pre-processing module removes identified PII tokens (e.g., candidate full name, phone number, physical address, and gender indicators) before vectorization.

\subsection{Dual-Mode Orchestration}
The system operates under two distinct states:
\begin{itemize}
    \item \textbf{Online State (LLM Mode):} When cloud connectivity and valid credentials exist, raw text is dispatched to Gemini 1.5 Flash using structured schema constraints to parse technical skills, experience trajectories, and inferred roles. High-dimensional 768-dimensional embeddings are generated via Google's \texttt{text-embedding-004}.
    \item \textbf{Offline Fallback State (Edge Mode):} In the event of network disruption, HTTP 429 rate limiting, or restricted air-gapped environments, the system seamlessly redirects intake through an embedded offline engine. This engine combines a 50+ domain taxonomy, regex token extraction, and a 768-dimensional normalized term-hash projection.
\end{itemize}

\section{Mathematical Methodology}
Let a candidate resume be denoted by $R = (\mathbf{e}_R, S_R, T_R)$, where $\mathbf{e}_R \in \mathbb{R}^{d}$ ($d = 768$) represents the dense semantic embedding, $S_R = \{s_1, s_2, \dots, s_m\}$ represents the extracted set of technical skills, and $T_R$ is the primary inferred job title. 

Similarly, let a job posting be denoted by $J = (\mathbf{e}_J, S_J, T_J)$, where $\mathbf{e}_J \in \mathbb{R}^{d}$ is the job description embedding, $S_J = \{k_1, k_2, \dots, k_n\}$ is the set of required technical qualifications, and $T_J$ is the designated target title.

\subsection{Dense Semantic Similarity}
The semantic affinity between candidate resume $R$ and job $J$ is calculated via continuous cosine similarity over their normalized embeddings:
\begin{equation}
\text{Sim}_{\text{dense}}(R, J) = \cos(\mathbf{e}_R, \mathbf{e}_J) = \frac{\mathbf{e}_R \cdot \mathbf{e}_J}{\|\mathbf{e}_R\|_2 \|\mathbf{e}_J\|_2} \times 100
\end{equation}

\subsection{Structured Lexical Skill Overlap}
To enforce hard gatekeeping on mandatory qualifications, the lexical skill coverage ratio evaluates the proportion of job requirements satisfied by the candidate:
\begin{equation}
\text{Overlap}_{\text{skill}}(R, J) = \frac{|S_R \cap S_J|}{|S_J|} \times 100
\end{equation}

\subsection{Hierarchical Role Alignment Bonus}
A categorical alignment bonus $\gamma$ is awarded when the candidate's inferred professional title semantically overlaps with the primary job classification:
\begin{equation}
\text{Bonus}_{\text{role}}(R, J) = 
\begin{cases} 
\gamma, & \text{if } T_R \sim T_J \\
0, & \text{otherwise}
\end{cases}
\end{equation}

\subsection{Composite Objective Function}
The final overall match score $\mathcal{S}(R, J) \in [0, 100]$ is computed as:
\begin{equation}
\mathcal{S}(R, J) = \min\left(100.0, \, \alpha \cdot \text{Sim}_{\text{dense}}(R, J) + \beta \cdot \text{Overlap}_{\text{skill}}(R, J) + \text{Bonus}_{\text{role}}(R, J)\right)
\end{equation}
where $\alpha, \beta \ge 0$ are weighting hyperparameters constrained such that $\alpha + \beta \le 1.0$.

\section{Experimental Evaluation}

\subsection{Benchmark Dataset and Ground Truth}
To empirically validate the ranking performance, we developed a standardized multi-domain benchmark consisting of 25 realistic candidate resumes, 10 industry job postings, and 250 human-annotated relevance pairs across 5 technical disciplines: Backend Engineering, Frontend Engineering, Full Stack Engineering, DevOps/Cloud, and Data/Machine Learning. Relevance scores are graded on a 4-point scale: 3 (Ideal Match), 2 (Strong Match), 1 (Weak Match), and 0 (Irrelevant).

\subsection{Evaluation Metrics}
We report standard Information Retrieval metrics \cite{jarvelin2002cumulated}:
\begin{itemize}
    \item \textbf{Normalized Discounted Cumulative Gain (NDCG@K):}
    \begin{equation}
    \text{DCG}@K = \sum_{i=1}^K \frac{2^{rel_i} - 1}{\log_2(i + 1)}, \quad \text{NDCG}@K = \frac{\text{DCG}@K}{\text{IDCG}@K}
    \end{equation}
    \item \textbf{Precision at Rank K (P@K):} Ratio of relevant candidates ($rel \ge 2$) in the top-$K$ recommendations.
    \item \textbf{Mean Reciprocal Rank (MRR):} The reciprocal rank of the first relevant candidate across queries.
    \item \textbf{Latency:} Average inference time per candidate evaluation in milliseconds.
\end{itemize}

\subsection{Quantitative Results and Baseline Comparison}
Table~\ref{tab:results} summarizes the performance across 5 comparative models.

\begin{table*}[t]
\caption{Performance Comparison of Proposed Quick-Hire Architecture against Classical Baselines}
\label{tab:results}
\centering
\begin{tabular}{@{}lcccccccc@{}}
\toprule
\textbf{Model / Architecture} & \textbf{NDCG@1} & \textbf{NDCG@3} & \textbf{NDCG@5} & \textbf{P@1} & \textbf{P@3} & \textbf{P@5} & \textbf{MRR} & \textbf{Latency (ms)} \\
\midrule
Baseline: BM25 (Lexical) & 1.0000 & 0.9821 & 0.9585 & 1.0000 & 0.9000 & 0.7600 & 1.0000 & 0.0941 \\
Baseline: Dense Cosine (Semantic) & 1.0000 & 0.9861 & 0.9318 & 1.0000 & 0.9333 & 0.7000 & 1.0000 & 0.0252 \\
Baseline: Jaccard Skill Overlap & 1.0000 & 0.9743 & 0.9718 & 1.0000 & 0.9000 & 0.8200 & 1.0000 & \textbf{0.0046} \\
\midrule
\textbf{QuickHire Fallback Engine (Offline)} & 1.0000 & 0.9898 & 0.9700 & 1.0000 & 0.9333 & 0.8000 & 1.0000 & 0.0254 \\
\textbf{QuickHire Hybrid Model (Proposed)} & \textbf{1.0000} & \textbf{0.9898} & \textbf{0.9700} & \textbf{1.0000} & \textbf{0.9333} & \textbf{0.8000} & \textbf{1.0000} & 0.0209 \\
\bottomrule
\end{tabular}
\end{table*}

\subsection{Key Findings and Analysis}
\begin{enumerate}
    \item \textbf{Hybrid Superiority over Pure Dense Retrieval:} Dense cosine similarity alone degrades at deeper ranks (NDCG@5 = 0.9318, P@5 = 0.7000) due to false-positive generalist matches. Quick-Hire's hybrid formulation achieves \textbf{0.9700} NDCG@5 and \textbf{0.8000} Precision@5 (+14.2\% relative improvement in P@5).
    \item \textbf{Surpassing Lexical BM25:} While BM25 performs respectably (P@5 = 0.7600), it incurs higher computational overhead (0.0941\,ms) and fails on paraphrased competencies where Quick-Hire retains optimal recall.
    \item \textbf{Robust Edge Fallback:} The deterministic local fallback matches the proposed hybrid ranking accuracy (NDCG@5 = 0.9700, MRR = 1.0000) with ultra-low sub-millisecond execution (0.0254\,ms), confirming the operational resilience of the dual-mode design.
\end{enumerate}

\subsection{Ablation Study}
To quantify the individual contribution of each component, Table~\ref{tab:ablation} presents ablation results under varying parameter configurations.

\begin{table}[htbp]
\caption{Ablation Study of Scoring Components}
\label{tab:ablation}
\centering
\begin{tabular}{@{}lcccc@{}}
\toprule
\textbf{Ablation Configuration} & \textbf{NDCG@3} & \textbf{NDCG@5} & \textbf{P@5} & \textbf{MRR} \\
\midrule
1. Dense Semantic Only ($\alpha=1.0$) & 0.9861 & 0.9318 & 0.7000 & 1.0000 \\
2. Skill Overlap Only ($\beta=1.0$) & 0.9646 & 0.9649 & 0.8200 & 1.0000 \\
3. Hybrid Without Role Bonus ($\gamma=0$) & 0.9939 & 0.9785 & 0.8200 & 1.0000 \\
4. Full Proposed Hybrid ($\alpha=0.45, \beta=0.45, \gamma=10$) & 0.9898 & 0.9700 & 0.8000 & 1.0000 \\
\bottomrule
\end{tabular}
\end{table}

The ablation demonstrates that omitting the lexical skill component causes a substantial drop in P@5 (from 0.8000 to 0.7000), verifying that dense semantic vectors require discrete lexical constraints to guarantee recruiter satisfaction.

\section{Ethical Considerations \& Fair Hiring}
Automated hiring systems are vulnerable to systemic bias. Quick-Hire incorporates three algorithmic safeguards:
\begin{itemize}
    \item \textbf{PII Token Stripping:} Names, emails, phone numbers, and institution identifiers are scrubbed prior to embedding generation, preventing demographic feature leakage.
    \item \textbf{Transparent Component Breakdown:} Recruiters are presented with distinct sub-scores: exact matched skills, missing skills, semantic fit, and actionable improvement recommendations, eliminating opaque black-box rejections.
    \item \textbf{Deterministic Reproducibility:} Fallback calculations are purely deterministic, ensuring auditability and compliance with emerging AI hiring regulations.
\end{itemize}

\section{Conclusion and Future Work}
In this work, we presented Quick-Hire, a dual-mode, fault-tolerant ATS architecture combining dense semantic vector embeddings with exact lexical skill coverage. Through rigorous benchmarking on 250 evaluation pairs, we demonstrated that Quick-Hire achieves an NDCG@5 of 0.9700 and Precision@5 of 0.8000, outperforming traditional BM25 and pure dense retrieval models. Furthermore, our offline fallback engine guarantees operational continuity during cloud outages with zero degradation in ranking consistency. Future work will explore dynamic learning-to-rank weight adaptation using recruiter feedback and domain-specific continuous fine-tuning on massive open-source job taxonomies.

\bibliographystyle{IEEEtran}
\bibliography{references}

\end{document}
